% python_api.tex — direct object-level Python API usage (GitTree / GitRepo)

\chapter{Direct Python Object API}
\label{chap:python-api}

This chapter shows how to manipulate \cgspkg{} objects directly (without the
CLI wrapper), including loading a \texttt{.gts} snapshot, reconstructing the
runtime registry, mirroring identity into \texttt{GitTree}/\texttt{GitRepo},
and propagating a tag target.

\section{From Empty Registry to READY via \texttt{.gts}}

\begin{lstlisting}[language=Python]
from ComplexGitSync import GitRepo, GitTree, RefKind, TreeLifecycleState
from ComplexGitSync.orchestre import GtsDocument, build_registry_from_gts_document
from ComplexGitSync.git_tree import DependencyTreeRegistry

# 1) Empty registry lifecycle
empty_registry = DependencyTreeRegistry()
assert empty_registry.recompute_tree_state() == TreeLifecycleState.UNLOADED

# 2) Load a runtime snapshot (.gts) and rebuild the dependency registry
gts_doc = GtsDocument.from_toml("examples/cawaqsviz_snapshot.gts")
registry = build_registry_from_gts_document(gts_doc)
assert registry.lifecycle_state == TreeLifecycleState.READY

# 3) Rebuild GitTree/GitRepo identity objects from discovered registry entries
tree = GitTree()
for entry in registry.values():
    tree.add_repo(
        GitRepo(
            project_owner_name=entry.project_owner_name or "owner",
            project_name=entry.project_name or entry.name,
            gitprovider=entry.gitprovider,
            access_protocol=entry.access_protocol,
            commit_sha=entry.commit_sha,
        )
    )
\end{lstlisting}

\section{Object-Level Tag Propagation}

\begin{lstlisting}[language=Python]
# Set the same tag target across all registry entries (parent -> leaves)
tree.propagate_tag(registry, "v1.2.3")
for entry in registry.values():
    assert entry.target_ref_kind == RefKind.TAG
    assert entry.target_ref_name == "v1.2.3"
\end{lstlisting}

This object-level flow is the same state model used by the CLI and
\texttt{ComplexGitSyncClient}, but exposed directly for advanced automation and
tests.

\section{READY Methods in the Client API}

The same lifecycle model is exposed by \texttt{ComplexGitSyncClient}.  In
practice, the workflow is:

\begin{itemize}
  \item \texttt{clone\_cgs(...)} bootstraps a tree and must finish in
    \texttt{READY}.
  \item \texttt{checkout(...)}, \texttt{add(...)}, \texttt{commit(...)},
    \texttt{push(...)}, and \texttt{tag(...)} are \texttt{READY}-gated methods
    for the standard git workflow.
  \item \texttt{freeze\_state(...)} and \texttt{launch\_state(...)} are the
    internal dev-state equivalents of release operations.
  \item \texttt{freeze\_release(...)} and \texttt{launch\_release(...)} are the
    public release operations using the same lifecycle model.
\end{itemize}

\begin{lstlisting}[language=Python]
from ComplexGitSync import ComplexGitSyncClient

client = ComplexGitSyncClient()

# Bootstrap from .cgs -> READY
client.clone_cgs("examples/complexgitsync.cgs")

# READY methods
client.checkout("feature/my-branch")
client.add()
client.commit("feat: update tree")
client.push()
client.tag("v1.2.3")

# Internal dev state (not publicly exposed as a release)
client.freeze_state("dev-state-2026.05")
client.launch_state(".cgitsync/state/dev-state-2026.05.gts")

# Public release
client.freeze_release("release-2026.05")

# Relaunch from a frozen release snapshot -> READY
client.launch_release(".cgitsync/releases/release-2026.05.gts")
\end{lstlisting}
